\documentclass[12pt,a4paper]{article}
\usepackage[T2A]{fontenc}
\usepackage[utf8]{inputenc}
\usepackage[russian]{babel}
\usepackage{geometry}
\usepackage{enumitem}
\usepackage{listings}
\usepackage{xcolor}
\usepackage{hyperref}
\usepackage{longtable}
\usepackage{array}
\geometry{left=25mm,right=15mm,top=20mm,bottom=20mm}
\setlength{\parindent}{1.25cm}
\setlength{\parskip}{0pt}
\lstset{basicstyle=\ttfamily\small,frame=single,breaklines=true,
showstringspaces=false,columns=fullflexible,keepspaces=true}
\hypersetup{colorlinks=true,linkcolor=black,urlcolor=blue}
\begin{document}
\begin{titlepage}
\begin{center}
Федеральное государственное автономное образовательное учреждение\\
высшего образования\\[2mm]
\textbf{«Национальный исследовательский университет\\
``Высшая школа экономики''»}\\[3mm]
Факультет компьютерных наук\\
Департамент программной инженерии
\vfill
{\Large\bfseries Отчёт по индивидуальному заданию № 3}\\[4mm]
{\large по дисциплине «Архитектура компьютера и операционные системы»}\\[7mm]
{\Large\bfseries Вариант 14. «Задача об умных цветах»}\\[3mm]
{\large Имитационное моделирование с помощью взаимодействующих процессов}
\vfill
\begin{flushright}
Выполнила: Гончарова Екатерина\\
Группа: БПИ-253
\end{flushright}
\vfill
Москва, 2026
\end{center}
\end{titlepage}
\tableofcontents
\newpage

\section{Постановка задачи}
Требуется разработать консольную программу для Linux, моделирующую работу умных цветов и садовников. Цветы постепенно теряют влагу и при достижении заданного порога отправляют заявку на полив. Садовники выбирают заявки, добираются до соответствующего цветка и выполняют полив. Одну заявку может обслуживать только один садовник. Если цветок высох до критического уровня, он погибает; повторный полив в течение защитного интервала также может привести к гибели. Ход моделирования отображается в консоли и записывается в журнал.

В третьем задании участники моделируются \textbf{процессами}, а не потоками. Процессы имеют независимые адресные пространства, поэтому обмен состоянием модели требует специальных средств межпроцессного взаимодействия (IPC). Подготовлены четыре варианта с разными средствами синхронизации и передачи данных.

\section{Модель и входные параметры}
\subsection{Сущности модели}
\texttt{Flower} содержит идентификатор цветка, влажность, скорость высыхания, текущее состояние, данные о заявке, оставшееся время полива, восстановления и защиты, а также число завершённых поливов. \texttt{Gardener} содержит идентификатор, занятость, индекс обслуживаемого цветка, объём оставшейся работы и скорость. \texttt{Settings} хранит конфигурацию всего сада.

Используются состояния \texttt{NORMAL}, \texttt{WILTING}, \texttt{WATERING}, \texttt{RECOVERING}, \texttt{DRY} и \texttt{DEAD\_OVERWATERED}. Начальная влажность каждого цветка составляет 100. После достижения порога заявки создаётся запрос на полив; если влажность достигла порога засыхания, заявка снимается и цветок перестаёт обслуживаться.

\subsection{Параметры запуска}
\begin{longtable}{|p{0.35\textwidth}|p{0.55\textwidth}|}
\hline
Параметр & Значение и ограничение\\ \hline
\endhead
Число цветов & От 1 до 100\\ \hline
Скорости высыхания & По одному положительному целому числу на цветок\\ \hline
Порог заявки & Целое число от 1 до 100\\ \hline
Порог засыхания & Целое число от 0 до 99; строго меньше порога заявки\\ \hline
Длительность полива & Положительное целое число тактов\\ \hline
Восстановление & Неотрицательное число тактов\\ \hline
Защитный интервал & Неотрицательное число тактов\\ \hline
Число садовников & От 1 до 10\\ \hline
Скорость садовника & Положительное целое число\\ \hline
Правило выбора & \texttt{first} либо \texttt{driest}\\ \hline
\end{longtable}

Например, при пороге заявки 30 и пороге засыхания 10 цветок запрашивает полив при снижении влажности до 30, а при дальнейшем снижении до 10 погибает. Правило \texttt{first} выбирает первую доступную заявку; \texttt{driest} --- заявку от цветка с минимальной влажностью. Время модели ограничено 50 тактами. В реализации используется задержка 25 мс на такт; это задержка отображения и выполнения, а не реальное время роста растений.

\section{Общая структура программы}
Основная модель разделена на небольшие файлы: \texttt{flower\_logic.c} обновляет отдельный цветок, \texttt{gardener\_logic.c} управляет поливом, \texttt{request\_logic.c} выбирает заявки, \texttt{model.c} инициализирует структуры, \texttt{input.c} и \texttt{config\_parser.c} читают параметры, \texttt{statistics.c} рассчитывает итоги, \texttt{logger.c} записывает события. В каждой реализации изменяются прежде всего \texttt{simulation.c}, функции процессов и IPC-модуль. Такой подход позволяет сравнивать механизмы взаимодействия без изменения правил поведения цветов.

В первых двух вариантах запускается \texttt{smart\_flowers}, в третьем и четвёртом --- три исполняемых файла: \texttt{garden}, \texttt{flower}, \texttt{gardener}. Главный процесс или контроллер отвечает за создание ресурсов, контроль времени, остановку и вывод итогов.

\section{Вариант 1. Неименованные семафоры и разделяемая память}
\subsection{Создание процессов}
Основная программа считывает \texttt{Settings}, создаёт POSIX-объект разделяемой памяти через \texttt{shm\_open()}, задаёт его размер с помощью \texttt{ftruncate()} и отображает в адресное пространство вызовом \texttt{mmap()}. Затем родитель создаёт дочерние процессы через \texttt{fork()}: один процесс на каждый цветок и каждый садовник.

Поскольку после \texttt{fork()} обычные локальные переменные процессов независимы, состояния цветов, заявки и признаки остановки размещаются в \texttt{SharedModel}, доступной всем процессам.

\subsection{Синхронизация}
Внутри разделяемой памяти инициализируются неименованные POSIX-семафоры через \texttt{sem\_init()} с параметром \texttt{pshared=1}. Семафор модели обеспечивает взаимное исключение при проверке и изменении заявок. Отдельные семафоры уведомляют садовников о появлении работы. Проверка свободной заявки и её закрепление выполняются в одной критической секции; два садовника не могут одновременно принять одну заявку.

При завершении родитель останавливает детей, дожидается их через \texttt{waitpid()}, уничтожает семафоры и удаляет shared memory. Обработка \texttt{SIGINT} позволяет выполнить тот же порядок очистки при \texttt{Ctrl+C}.

\section{Вариант 2. Именованные семафоры и разделяемая память}
Процессы запускаются тем же способом --- родителем через \texttt{fork()}; память также создаётся через \texttt{shm\_open()} и \texttt{mmap()}. Отличается тип семафоров. Вместо \texttt{sem\_init()} используются \texttt{sem\_open()} и уникальные имена семафоров, связанные с запуском программы.

Процесс получает доступ к именованному семафору по его имени. Для завершения вызываются \texttt{sem\_close()} и \texttt{sem\_unlink()}. Это важно: закрытие дескриптора само по себе не удаляет имя объекта из системы. Модель цветов и правила назначения заявок в первом и втором вариантах одинаковы. Разница заключается именно в организации семафоров.

\section{Вариант 3. Независимые процессы и разделяемая память}
\subsection{Запуск}
Третий вариант не создаёт рабочие процессы через \texttt{fork()} внутри контроллера. Вместо этого отдельные исполняемые файлы \texttt{garden}, \texttt{flower} и \texttt{gardener} могут запускаться из разных терминалов. Контроллер \texttt{garden} создаёт именованные семафоры и разделяемую память, а работники подключаются к ним по имени сессии. Каждый работник получает роль и номер, начиная с 1.

Контроллер ждёт регистрации ожидаемого числа цветов и садовников. Для предотвращения повторного подключения одного работника проверяются идентификаторы и занятые позиции. Если не все процессы зарегистрировались в течение 15 секунд, запуск завершается с ошибкой и очисткой доступных ресурсов.

\subsection{Обмен данными}
Все участники читают и изменяют одну общую структуру в POSIX shared memory. Именованный семафор модели ограничивает доступ к этой структуре, чтобы состояние заявки и назначение садовника изменялись согласованно. Уведомления о новых заявках передаются через отдельные именованные семафоры садовников.

Контроллер не является родителем работников и поэтому не может обычным \texttt{waitpid()} дождаться этих независимых процессов. Он проверяет их регистрацию и доступность, устанавливает общий признак остановки и удаляет собственные IPC-ресурсы при завершении сессии.

\section{Вариант 4. Независимые процессы и очереди сообщений}
\subsection{Архитектура}
Четвёртый вариант также состоит из независимо запускаемых программ \texttt{garden}, \texttt{flower} и \texttt{gardener}, но \textbf{состояние модели не передаётся через общую память}. Контроллер хранит основной экземпляр модели в собственной памяти. Для обмена с работниками используются POSIX-очереди сообщений, создаваемые через \texttt{mq\_open()}.

Для сессии создаются очередь событий, через которую работники передают регистрацию и результаты, и отдельные очереди заданий для работников. Контроллер передаёт снимок состояния, роль и номер такта. Получив сообщение, процесс цветка или садовника применяет свою часть логики и отправляет результат обратно. Контроллер проверяет номер отправителя и такта и вносит изменение в основную модель.

\subsection{Семафоры и ограничения}
Именованный семафор используется для уведомления контроллера о поступившем событии. Дополнительные именованные семафоры не допускают двойной регистрации одного номера работника. Последовательное применение результатов садовников предотвращает конфликтующие изменения состояния заявки.

В очереди каждого работника достаточно одного ожидающего задания: следующий шаг отправляется только после обработки предыдущего. Это также уменьшает объём памяти IPC и позволяет запускать конфигурации с большим количеством рабочих процессов без избыточного резервирования очередей. При завершении контроллер отправляет сообщения остановки и удаляет имена очередей через \texttt{mq\_unlink()}, а семафоров --- через \texttt{sem\_unlink()}.

\section{Пользовательский интерфейс и запуск}
Все команды выполняются в Linux из каталога соответствующего варианта. Компиляция проводится с помощью \texttt{make}, автоматические сценарии --- с помощью \texttt{make check}. Для первых двух вариантов поддерживаются аргумент с файлом конфигурации и ручной ввод параметров без файла:
\begin{lstlisting}
make
./build/smart_flowers tests/balanced.cfg
./build/smart_flowers
\end{lstlisting}

В вариантах 3 и 4 можно запустить контроллер с файлом и именем сессии или с одним именем для ручного ввода:
\begin{lstlisting}
make
./build/garden tests/one_flower.cfg demo
./build/garden demo
\end{lstlisting}

Во втором случае контроллер последовательно запрашивает число цветов, индивидуальные скорости высыхания, пороги, длительности, число садовников, скорость и правило выбора заявки. Некорректное число можно ввести повторно. После настройки остальные процессы подключаются из других терминалов:
\begin{lstlisting}
./build/flower demo 1
./build/gardener demo 1
\end{lstlisting}

Для конфигураций с несколькими участниками требуется запустить экземпляр программы для каждого номера. Сценарий \texttt{scripts/run\_demo.sh} автоматически определяет количество работников из конфигурационного файла и запускает их в отдельных фоновых процессах:
\begin{lstlisting}
bash scripts/run_demo.sh tests/balanced.cfg
\end{lstlisting}

При таком запуске имя сессии выбирается скриптом автоматически. Для ручного запуска нужно задать своё имя: от 1 до 24 латинских букв, цифр, символов подчёркивания или дефисов. Все ключевые события видны в консоли. Результаты сохраняются в \texttt{simulation.log} либо \texttt{simulation\_\textit{session}.log} в независимых вариантах. В штатной ситуации симуляция завершается после 50 тактов или когда все цветы стали неактивны. Нажатие \texttt{Ctrl+C} в терминале главной программы останавливает выполнение раньше.

\section{Проверки и результаты работы}
\subsection{Набор сценариев}
Для вариантов предусмотрены конфигурации, проверяющие работу нескольких садовников, стратегию \texttt{driest}, быстрое высыхание, засыхание в начале работы, длительное восстановление, отсутствие защитного интервала, перелив, одного работника, большое число процессов и крайние допустимые значения порогов. Всего используется 13 файлов конфигурации.

Дополнительно тестируются некорректные аргументы, недопустимые значения параметров, прерывание по \texttt{Ctrl+C}, корректное завершение работников и очистка IPC-ресурсов. В третьем и четвёртом вариантах отдельно проверяются независимое подключение, повторная регистрация и недопустимые номера процессов.

\subsection{Проверка интерактивного запуска}
Для вариантов 3 и 4 был отдельно проверен ручной ввод с одним цветком и одним садовником. Сначала в качестве порога засыхания было введено число, превышающее порог заявки; программа сообщила об ошибке и запросила значение повторно. После исправления параметра, запуска \texttt{flower} и \texttt{gardener} контроллер дождался обоих участников, выполнил симуляцию и вывел итоговую статистику. Обе проверки завершились успешно.

\subsection{Результаты тестирования}
Для каждого варианта использовались 13 конфигурационных файлов с различными входными параметрами. Проверки запускались командой \texttt{make check}. Результаты приведены в таблице.

\begin{longtable}{|p{0.60\textwidth}|p{0.25\textwidth}|}
\hline
Реализация & Результат \\ \hline
\endhead
Вариант 1 (неименованные семафоры) & 13 сценариев пройдены \\ \hline
Вариант 2 (именованные семафоры) & 13 сценариев пройдены \\ \hline
Вариант 3 (независимые процессы, разделяемая память) & 13 сценариев пройдены \\ \hline
Вариант 4 (независимые процессы, очереди сообщений) & 13 сценариев пройдены \\ \hline
\end{longtable}

Все четыре реализации успешно прошли предусмотренные тестовые сценарии.

Итоговый вывод программы содержит последний обработанный такт, число живых и погибших цветов, общее число поливов, а также состояние и влажность каждого цветка. Журнал позволяет проследить создание заявки, её принятие, движение садовника, завершение полива и переходы между состояниями. Конкретные числовые результаты зависят от выбранной конфигурации и порядка выполнения процессов, поэтому для воспроизводимого наблюдения целесообразно использовать конфигурации из каталога \texttt{tests/}.

\subsection{Пример результата работы}
Для примера, при запуске четвёртой реализации с \texttt{tests/one\_flower.cfg} один из запусков завершился на 50-м такте со следующими показателями: один живой цветок, ноль засохших, ноль погибших от перелива и семь завершённых поливов. В первых трёх вариантах число операций и промежуточные состояния могут отличаться, поскольку процессы обновляют модель асинхронно.

\section{Особенности и ограничения реализации}
В первых трёх вариантах общая разделяемая структура защищается семафором. Такой способ прост и позволяет использовать общий код модели, но при большом числе работников ограничивает параллельное изменение данных. В четвёртом варианте основные изменения последовательно применяет контроллер. Это устраняет конфликтующие записи, но каждый шаг обслуживания требует передачи сообщений и ожидания ответа.

Ограничение на число POSIX-очередей и объём выделяемой им памяти зависит от системы Linux. В четвёртом варианте персональные очереди имеют вместимость в одно задание, что соответствует последовательному протоколу обмена. При жёстком принудительном завершении \texttt{SIGKILL} нельзя гарантировать удаление всех именованных объектов: этот сигнал не обрабатывается программой. При нормальном завершении и \texttt{Ctrl+C} выполняются предусмотренные операции очистки.

\section{Заключение}
В рамках задания реализованы четыре способа организации взаимодействия процессов. Первые два варианта показывают различие между неименованными и именованными POSIX-семафорами при использовании общего родителя. Третий позволяет запускать участников независимо, подключая их к общей памяти по имени сессии. Четвёртый использует обмен сообщениями без разделяемой памяти модели. Поведение цветов и садовников во всех реализациях построено на одной предметной модели, поэтому различия относятся прежде всего к организации процессов, передаче данных и синхронизации.
\end{document}